\documentclass[letterpaper,10pt,onecolumn]{article}
\usepackage[spanish]{babel}
\usepackage[latin1]{inputenc}
\usepackage{amsfonts}
\usepackage{amsthm}
\usepackage{amsmath}
\usepackage{mathrsfs}
\usepackage{empheq}
\usepackage{enumitem}
\usepackage[pdftex]{color,graphicx}
\usepackage{hyperref}
\usepackage{listings}
\usepackage{calligra}
\usepackage{algpseudocode} 
\DeclareMathAlphabet{\mathcalligra}{T1}{calligra}{m}{n}
\DeclareFontShape{T1}{calligra}{m}{n}{<->s*[2.2]callig15}{}
\newcommand{\scripty}[1]{\ensuremath{\mathcalligra{#1}}}
\lstloadlanguages{[5.2]Mathematica}
\setlength{\oddsidemargin}{0cm}
\setlength{\textwidth}{490pt}
\setlength{\topmargin}{-40pt}
\addtolength{\hoffset}{-0.3cm}
\addtolength{\textheight}{4cm}

\begin{document}
\begin{center}

\includegraphics[width=490pt]{header.png}\\[0.5cm]

\textsc{\LARGE Taller 2 - F\'isica I (FISI-1018) - 2015-10}\\[0.5cm]

\textsc{\Large{Profesor: Jaime Forero}} \\[0.5cm]

\noindent\textsc{Ejercicios correspondiente a la clase complementaria
  de la semana del 2 de Febrero del 2015.}\\[0.5cm]
\end{center}

\noindent\textsc{Nota:} Los primeros cinco ejercicios deben ser
entregados {\bf al comienzo} de la clase complementaria. Los siguientes
cinco deben ser trabajados durante la complementaria. La numeraci\'on
hace referencia al texto gu\'ia: \textit{F\'isica Universitaria Volumen
  1 (Sears-Semansky)}, decimotercera edici\'on, Pearson.

\begin{enumerate}
\item Ejercicio 1.66 (Equilibrio de fuerzas).
\item Ejercicio 1.91 (Cubo).
\item Ejercicio 2.8 (Ave volando hacia el oriente).
\item Ejercicio 2.10 (Profesor caminando).
\item Ejercicio 2.29 (Lanzamiento del transbordador espacial).
\item Ejercicio 2.35 (Salto de una pulga).
\item Ejercicio 2.42 (Ladrillo cayendo de un edificio).
\item Ejercicio 2.54 (Salto volador de una pulga).
\item Probema 2.64 (Tren subterr\'aneo).
\end{enumerate}

\end{document}
